\begin{tikzpicture}[figure picture,x=10mm,y=10mm,font=\figurefont\fontsize{8}{10}\selectfont]
\path[use as bounding box] (0,0) rectangle (16.4,4.4);
% 左：度量
\node[figure heading,anchor=west] at (0.2,4.05) {度量 $d(x,y)$：输入是两个点，输出它们之间的距离};
\draw[FigureGrid,line width=.8pt] (0.6,1.5)--(7.4,1.5);
\fill[FigureInk] (2.4,1.5) circle[radius=1.1mm];
\node[below] at (2.4,1.35) {$x$};
\fill[FigureInk] (5.6,1.5) circle[radius=1.1mm];
\node[below] at (5.6,1.35) {$y$};
\draw[FigureRed,line width=1.1pt,{Stealth[length=2mm]}-{Stealth[length=2mm]}] (2.4,2.45)--(5.6,2.45);
\node[above,FigureRed] at (4.0,2.5) {$d(x,y)$};
\draw[FigureGrid,dashed,line width=.5pt] (2.4,1.5)--(2.4,2.4);
\draw[FigureGrid,dashed,line width=.5pt] (5.6,1.5)--(5.6,2.4);
\node[figure note,anchor=west,align=left] at (0.6,.55) {只问“两点隔多远”，\\不关心任何一块区域有多重};
% 分隔线
\draw[FigureGrid,line width=.5pt,dashed] (8.0,.35)--(8.0,4.15);
% 右：测度
\node[figure heading,anchor=west] at (8.3,4.05) {测度 $\mu(A)$：输入是一个集合，输出这一整块的份量};
\draw[FigureGrid,line width=.8pt] (8.7,1.5)--(15.8,1.5);
\fill[FigureBlueFill] (9.5,1.36) rectangle (13.6,1.64);
\draw[FigureBlue,line width=1.1pt] (9.5,1.36) rectangle (13.6,1.64);
\node at (11.55,1.5) {$A$};
\foreach \px in {9.9,10.6,11.2,11.9,12.6,13.2} \fill[FigureBlue] (\px,1.5) circle[radius=.55mm];
\draw[FigureBlue,{Stealth[length=1.6mm]}-] (11.55,1.32)--(11.55,.95);
\node[below,FigureBlue] at (11.55,.8) {$\mu(A)$};
\node[figure note,anchor=west,align=left] at (8.7,.35) {只问“这一块有多重”，\\完全不读两点之间的距离};
\end{tikzpicture}
